% Part: second-order-logic
% Chapter: sol-and-sets
% Section: introduction

\documentclass[../../../include/open-logic-section]{subfiles}

\begin{document}

\olfileid[tr]{sol}{set}{int}

\olsection{Giriş}

İkinci dereceden mantık fonksiyonların yanı sıra tanım kümesinin alt kümeleri
üzerinde de niceleme yapabildiğinden, küme teorisinin en azından bir
bölümünün ikinci dereceden mantıkta yürütülebilmesi beklenir. ``Yürütmek''
ile, küme kuramsal özellik ve ifadelerin ikinci dereceden mantıkta, kümeler
için özel ve mantık dışı bir söz varlığı (örneğin küme teorisinin üyelik
!!{predicate}{}i) olmadan ifade edilebilmesini kastediyoruz. Örneğin kümelerin
birleşim ve kesişimlerini ve alt küme bağıntısını tanımlayabilir; ayrıca
kümelerin büyüklüklerini karşılaştırıp Cantor Teoremi gibi sonuçları ifade
edebiliriz.

\end{document}
